\chapter{Conclusioni}
\label{chap:conclusioni}

\section{Considerazioni finali}
\noindent Il progetto di stage, della durata complessiva di trecento ore, si è concluso con successo e ha permesso di sviluppare e integrare il modulo completo e funzionale per la gestione delle note spesa aziendali all'interno dell'applicazione \gls{ergteam}. L'applicativo ha raggiunto l'obiettivo primario di modernizzare e snellire i processi amministrativi, riducendo l'uso del cartaceo e automatizzando l'inserimento dei dati grazie all'intelligenza artificiale. \\
Durante il percorso di sviluppo sono stati affrontati e risolti vari problemi tecnici e progettuali, mantenendo sempre una forte attenzione alla qualità del codice, alla reattività dell'interfaccia e all'esperienza utente. Particolare rilevanza è stata data all'implementazione di un'architettura "offline-first", garantendo la totale continuità operativa del lavoratore anche in condizioni di assenza di connettività di rete, e alla gestione sicura e crittografata dei dati sensibili in transito.
Il sistema finale è stato progettato per essere robusto, scalabile e facilmente manutenibile nel tempo, grazie alla rigorosa applicazione del pattern architetturale \gls{mvvm}. Questa scelta ha consentito una gestione modulare delle diverse componenti, separando in modo netto la logica di business dal livello di presentazione e consegnando all'azienda un prodotto moderno, efficiente e pronto per l'utilizzo in produzione.



\section{Raggiungimento degli obiettivi}
\noindent Al termine delle fasi di sviluppo, collaudo e validazione con il tutor aziendale, è stato verificato lo stato di implementazione di ogni singola funzionalità rispetto a quanto pattuito in fase di analisi nel Capitolo \ref{chap:analisi-requisiti}. Tutti gli obiettivi prefissati sono stati raggiunti con successo, come evidenziato nelle tabelle \ref{tab:conclusione_requisiti_funzionali}, \ref{tab:conclusione_requisiti_qualitativi} e \ref{tab:conclusione_requisiti_vincolo}, che riportano lo stato di soddisfacimento dei requisiti funzionali, qualitativi e di vincolo.

\begin{center}
    \rowcolors{1}{}{tableGray}
    \begin{longtable}{|p{2.15cm}|p{7.75cm}|p{2.35cm}|}
    \hline
    \multicolumn{1}{|c|}{\textbf{Requisito}} & \multicolumn{1}{c|}{\textbf{Descrizione}} & \multicolumn{1}{c|}{\textbf{Stato}}\\
    \hline 
    \endfirsthead
    \rowcolor{white}
    \multicolumn{3}{c}{{\bfseries \tablename\ \thetable{} -- Continuo della tabella}}\\
    \hline
    \multicolumn{1}{|c|}{\textbf{Requisito}} & \multicolumn{1}{c|}{\textbf{Descrizione}} & \multicolumn{1}{c|}{\textbf{Stato}}\\
    \hline 
    \endhead
    \hline
    \rowcolor{white}
    \multicolumn{3}{|r|}{{Continua nella prossima pagina...}}\\
    \hline
    \endfoot
    \endlastfoot
    
    RFO-1 & Visualizzazione lista resoconti mensili. & Soddisfatto \\ \hline
    RFO-2 & Possibilità di filtrare i resoconti mensili per anno. & Soddisfatto \\ \hline
    RFO-3 & Inserimento di una nuova nota spesa. & Soddisfatto \\ \hline
    RFO-4 & Validazione dei dati inseriti e annullamento salvataggio. & Soddisfatto \\ \hline
    RFO-5 & Inserimento descrizione nota spesa. & Soddisfatto \\ \hline
    RFO-6 & Possibilità di allegare una fotografia per estrazione dati con \gls{ia}. & Soddisfatto \\ \hline
    RFO-7 & Visualizzazione e modifica dati estratti. & Soddisfatto \\ \hline
    RFO-8 & Rimozione foto precedentemente allegata. & Soddisfatto \\ \hline
    RFO-9 & Aggiunta transazioni multiple. & Soddisfatto \\ \hline
    RFO-10 & Gestione del salvataggio locale e offline. & Soddisfatto \\ \hline
    RFO-11 & Sincronizzazione in blocco o singolarmente delle note spesa. & Soddisfatto \\ \hline
    RFO-12 & Visualizzazione lista note spesa di un resoconto mensile. & Soddisfatto \\ \hline
    RFO-13 & Eliminazione nota spesa esistente. & Soddisfatto \\ \hline
    RFO-14 & Impedimento modifica note spesa antecedenti alla data minima consentita. & Soddisfatto \\ \hline
    
    RFD-1 & Gestione dei conflitti tra dati inseriti manualmente e dati estratti dall'\gls{ia}. & Soddisfatto \\ \hline
    RFD-2 & Possibilità di rimozione spesa salvata temporaneamente tramite inserimento multiplo. & Soddisfatto \\ \hline
    RFD-3 & Promemoria presenza note locali non sincronizzate. & Soddisfatto \\ \hline
    RFD-4 & Possibilità di modificare una nota spesa. & Soddisfatto \\ \hline
    RFD-5 & Salvataggio modifiche a una nota spesa esistente. & Soddisfatto \\ \hline
    \hiderowcolors
    \caption{Tracciamento dello stato dei requisiti funzionali.}
    \label{tab:conclusione_requisiti_funzionali}
    \end{longtable}
\end{center}

\begin{center}
    \rowcolors{1}{}{tableGray}
    \begin{longtable}{|p{2.15cm}|p{7.75cm}|p{2.35cm}|}
    \hline
    \multicolumn{1}{|c|}{\textbf{Requisito}} & \multicolumn{1}{c|}{\textbf{Descrizione}} & \multicolumn{1}{c|}{\textbf{Stato}}\\
    \hline 
    \endfirsthead
    \rowcolor{white}
    \multicolumn{3}{c}{{\bfseries \tablename\ \thetable{} -- Continuo della tabella}}\\
    \hline
    \multicolumn{1}{|c|}{\textbf{Requisito}} & \multicolumn{1}{c|}{\textbf{Descrizione}} & \multicolumn{1}{c|}{\textbf{Stato}}\\
    \hline 
    \endhead
    \hline
    \rowcolor{white}
    \multicolumn{3}{|r|}{{Continua nella prossima pagina...}}\\
    \hline
    \endfoot
    \endlastfoot
    
    RQO-1 & Comunicazione chiara dello stato delle operazioni tramite messaggi bloccanti o temporanei. & Soddisfatto \\ \hline
    RQO-2 & Architettura del codice \textit{frontend} rispetta il pattern \gls{mvvm}. & Soddisfatto \\ \hline
    RQO-3 & Test di accettazione per validare le funzionalità. & Soddisfatto \\ \hline
    RQO-4 & Realizzazione manuale utente per l'utilizzo dell'applicazione. & Soddisfatto \\ \hline
    RQD-1 & Commento adeguato del codice per facilitarne la manutenzione. & Soddisfatto \\ \hline
    \hiderowcolors
    \caption{Tracciamento dello stato dei requisiti qualitativi.}
    \label{tab:conclusione_requisiti_qualitativi}
    \end{longtable}
\end{center}

\begin{center}
    \rowcolors{1}{}{tableGray}
    \begin{longtable}{|p{2.15cm}|p{7.75cm}|p{2.35cm}|}
    \hline
    \multicolumn{1}{|c|}{\textbf{Requisito}} & \multicolumn{1}{c|}{\textbf{Descrizione}} & \multicolumn{1}{c|}{\textbf{Stato}}\\
    \hline 
    \endfirsthead
    \rowcolor{white}
    \multicolumn{3}{c}{{\bfseries \tablename\ \thetable{} -- Continuo della tabella}}\\
    \hline
    \multicolumn{1}{|c|}{\textbf{Requisito}} & \multicolumn{1}{c|}{\textbf{Descrizione}} & \multicolumn{1}{c|}{\textbf{Stato}}\\
    \hline 
    \endhead
    \hline
    \rowcolor{white}
    \multicolumn{3}{|r|}{{Continua nella prossima pagina...}}\\
    \hline
    \endfoot
    \endlastfoot
    
    RVO-1 & Sviluppo in \gls{csharpg} utilizzando il framework \textit{cross-platform} \gls{dotnetmauig}. & Soddisfatto \\ \hline
    RVO-2 & Interfaccia grafica realizzata con \gls{xamlg}. & Soddisfatto \\ \hline
    RVO-3 & Nuovi endpoint REST nel \textit{backend}. & Soddisfatto \\ \hline
    RVO-4 & L'applicazione deve funzionare su dispositivi Android. & Soddisfatto \\ \hline
    RVO-5 & L'applicazione deve funzionare su dispositivi iOS. & Soddisfatto \\ \hline
    RVO-6 & Dati delle note spesa devono integrarsi con l'\gls{erp} aziendale. & Soddisfatto \\ \hline
    RVO-7 & Interrogazione dell'\gls{api} di Azure Content Understanding per estrazione dati. & Soddisfatto \\ \hline
    \hiderowcolors
    \caption{Tracciamento dello stato dei requisiti di vincolo.}
    \label{tab:conclusione_requisiti_vincolo}
    \end{longtable}
\end{center}


\section{Conoscenze acquisite}
\noindent Durante lo sviluppo del progetto, ho avuto l'opportunità di acquisire e consolidare competenze avanzate in diverse aree dell'ingegneria del software. Dal punto di vista tecnico, l'apprendimento del framework \gls{dotnetmauig} e dei linguaggi \gls{csharpg} e \gls{xamlg} mi ha permesso di comprendere a fondo come progettare e realizzare un'applicazione mobile cross-platform moderna, applicando rigorosamente il pattern architetturale \gls{mvvm}. L'integrazione del database locale \gls{sqliteg} per garantire un'architettura "offline-first" e l'utilizzo delle \gls{api} di \gls{azurecontentunderstandingg} per l'estrazione dei dati tramite intelligenza artificiale hanno rappresentato sfide altamente stimolanti che hanno arricchito notevolmente il mio bagaglio tecnologico.
Oltre alle competenze puramente informatiche, l'inserimento in un contesto aziendale dinamico mi ha permesso di migliorare in modo significativo le mie competenze trasversali. Ho avuto modo di affinare le mie capacità di problem solving affrontando ostacoli reali, imparando a pianificare le attività di sviluppo rispettando le scadenze e comprendendo l'importanza del lavoro in team e della comunicazione continua con il tutor aziendale.


\section{Valutazione personale}
\noindent Questo stage ha rappresentato un'opportunità formativa preziosa, permettendomi di consolidare le conoscenze apprese durante il percorso accademico e di metterle in pratica in un contesto aziendale reale. La progettazione e lo sviluppo dell'applicazione mobile hanno richiesto un impegno costante, applicando praticamente le conoscenze acquisite per gestire l'intero ciclo di vita di un prodotto software, dall'analisi dei requisiti fino al collaudo finale. L'ambiente di Ergon Informatica si è rivelato estremamente accogliente e stimolante: la collaborazione continua con il tutor e altri colleghi ha arricchito ulteriormente l'esperienza, insegnandomi l'importanza della comunicazione e della pianificazione nello sviluppo software. Nel complesso, sono molto soddisfatto dei risultati ottenuti e sono fiducioso che le competenze acquisite avranno un impatto altamente positivo sul mio futuro professionale.


\newpage